\documentclass[10pt,a4paper]{amsart}
\usepackage[margin=19mm]{geometry}
\usepackage{amsmath,amssymb,mathtools}
\usepackage[scr=rsfs,cal=euler]{mathalfa}
\usepackage{xfrac}
\usepackage[unicode,hidelinks,bookmarks=false]{hyperref}
\newcommand{\ie}{i.e.\ }
\newcommand{\cf}{cf.\ }
\newcommand{\resp}{respectively\ }
\newcommand{\Subsection}{\subsection}
\providecommand{\nobreakdash}{\nobreak\hbox{-}}
\setlength{\emergencystretch}{2em}
\multlinegap0pt
\usepackage{amssymb}
\usepackage{mathtools}
\newtheorem{defn}{}
\newtheorem{prop}{}
\newtheorem{lem}{}
\DeclareMathOperator{\IM}{Im}
\DeclareMathOperator{\RE}{Re}
\def\T{{\mathbb T}}% torus
\def\d{{\partial}}
\title{The time-relaxation limit for weak solutions to the quantum hydrodynamics system}
\author{Paolo Antonelli, Pierangelo Marcati, Hao Zheng}
\date{Source: arXiv:2412.16800v2 (CC BY 4.0)}
\begin{document}
\maketitle

\noindent\emph{Source and licence.} The time-relaxation limit for weak solutions to the quantum hydrodynamics system. Version arXiv:2412.16800v2 (CC BY 4.0), distributed by arXiv under the Creative Commons Attribution 4.0 International licence, http://creativecommons.org/licenses/by/4.0/, as recorded in the arXiv licence metadata for this exact version and shown on the abstract page https://arxiv.org/abs/2412.16800v2. Full licence notice: \texttt{licenses/arxiv-2412.16800-CC-BY-4.0.txt}.\par\smallskip

\noindent\textit{Editorial scope.} Counted unit: the article's prop eq:ini\_1 (source0.tex lines 665--671). The article's complete original proof of it is delivered with it (source0.tex lines 673--835, one \texttt{proof} environment of 163 lines). No mathematics is written, repaired or re-derived: the statement and the proof are the article's own lines, in the article's own notation, and no retained line is substituted or re-flowed.  Retained uncounted as the source-local context the proof needs: lines 97--103; lines 194--196; lines 548--553; lines 564--576; lines 657--662.  Nothing was omitted from any retained span, so the package carries no omission record.  No retained line contains a citation, so the wrapper carries no bibliography and no bibliography entry.  No retained line contains a figure environment, an image-inclusion command or drawing code, so no image is delivered and the package declares no asset dependency.  Editorial formatting only, in addition: an AMS \texttt{amsart} preamble that restates the article's own declarations and macros used by the retained lines, namely the environments \texttt{defn}, \texttt{prop}, \texttt{lem} on separate counters, together with the standard packages those lines need.  The retained environments print in this document as Definition 1, Proposition 1, Lemma 1 in order of appearance, whereas the article prints the same items with the numbering of its own full document.  The complete native source of the article as distributed in the arXiv e-print for this exact version is bundled unchanged as \texttt{sources/170541/source0.tex}.
\begin{equation}\label{eq:QHD}
\begin{cases}
\d_{t}\rho+\d_xJ=0\\
\d_{t}J+\d_x\left(\frac{J^2}{\rho}\right)+\d_x p(\rho)+\rho\d_x V=\frac{1}{2}\rho\d_x\left(\frac{\d_x^2\sqrt\rho}{\sqrt\rho}\right)-\frac{1}{\tau}J\\
-\d_x^2V=\rho-\mathcal{C}(x),\quad (\rho,J)(0,x)=(\rho_0,J_0)(x),
\end{cases}
\end{equation}

\begin{equation}\label{eq:chem}
\mu=-\frac{\d_x^2\sqrt\rho}{2\sqrt\rho}+\frac12\frac{J^2}{\rho^2}+f'(\rho)+V
\end{equation}

\begin{equation}\label{eq:gradS}
\begin{cases}
\d_xS =v,\quad\d_t S=-\mu-\frac{1}{\tau}S \\
S(0,0)=S_*.
\end{cases}
\end{equation}

\begin{defn}[Phase function]\label{def:phase}
Let $(\rho,v)$ be a pair a solution of \eqref{eq:QHD} on $[0,T)\times\T$ such that the velocity $v$ and the chemical potential $\mu$ given by \eqref{eq:chem} satisfy, 
\[
v\in L^\infty_tL^1_x,\quad\mu\in \mathcal{D}([0,T)\times\T),
\]
and $\int_\T v dx=0$. Then we define the phase function $S$ to be the solution of the gradient equation \eqref{eq:gradS}, which can be explicitly written as
\begin{equation}\label{eq:defS}
\begin{aligned}
S(t,x)=&e^{-\frac{t}{\tau}} \left(S_*+\int_\T\int_0^{x_*}v_0(y)dydx_*\right)\\
&-\int_0^t \int_\T e^{\frac{s-t}{\tau}}\mu(s,x_*)dx_*ds+\int_\T\int_{x_*}^x v(t,y)dydx_*.
\end{aligned}
\end{equation}
\end{defn}

\begin{equation}\label{eq:cauchy_NLS}
\begin{cases}
i\d_t\psi+\frac12\d_x^2\psi=f'(|\psi|^2)\psi+\frac{1}{\tau}S\psi+V\psi\\
-\d_x^2 V=|\psi|^2-\mathcal{C}(x),\quad \psi(0)=\psi_0
\end{cases}
\end{equation}

\begin{prop}
Let us assume the initial data $\psi_0\in H^2(\T)$ satisfy $\inf_x|\psi_0|\geq \delta^\frac12 >0$ and 
\begin{equation}\label{eq:ini_1}
\|\d_x\psi_0\|_{L^2_x} \leq M_0^\frac12-\delta M_0^{-\frac12},
\end{equation}
where $M_0=\|\psi_0\|_{L^2_x}^2$. Then the Cauchy problem \eqref{eq:cauchy_NLS} has a unique solution $\psi\in L^\infty(0,T_*;H^2(\T))$ local in time, for some $T_*>0$, with $\inf_{t,x}|\psi|\geq \frac{\delta^\frac12}{2}$ .
\end{prop}

\begin{proof}
The proof follows a standard Picard iteration argument. The sequence $\{\psi_n\}_{n\geq 1}$ for \eqref{eq:cauchy_NLS} is defined iteratively by solving $\psi_n$ from the linear Schr\"odinger equation 
\begin{equation}\label{eq:QLS_n}
i\d_t\psi_n+\frac12\d_x^2\psi_n=f'(|\psi_{n-1}|^2)\psi_n+\frac{1}{\tau}S_{n-1}\psi_n+V_{n-1}\psi_n
\end{equation}
on $[0,T_*)\times \T$, for small $T_*$ to be determined later, with $\psi_{n-1}$, $S_{n-1}$ and $V_{n-1}$ given in the previous step. The Duhamel formula for \eqref{eq:QLS_n} is written as
\begin{equation}\label{eq:duhamel}
\begin{aligned}
\psi_n=e^{\frac{i}{2}t\d_x^2}\psi_0&-i\int_0^te^{\frac{i}{2}(t-s)\d_x^2}[f'(|\psi_{n-1}|^2)+V_{n-1}]\psi_n(s)ds\\
&-\frac{i}{\tau}\int_0^te^{\frac{i}{2}(t-s)\d_x^2}S_{n-1}\psi_n(s)ds,
\end{aligned}
\end{equation}
and the electric potential $V_{n-1}$ is provided by the Poisson equation, 
\[
-\d_x^2 V_{n-1}=|\psi_{n-1}|^2-\mathcal{C}(x)
\]
with $\int_\T V_{n-1}dx=0$. By Poincar\'e inequality, $V_{n-1}$ is bounded by
\[
\|V_{n-1}\|_{W^{1,p}_x}\leq C(\|\psi_{n-1}\|^2_{L^2_x}+\|\mathcal{C}\|_{L^1_x}),\quad 1\leq p\leq \infty.
\]
For the phase function, we first define $S_0=0$. For $n\geq 1$, we want to give s suitable definition of $S_n$ by adapting the idea of Definition \ref{def:phase} for $\psi_n$ as follows. As before, we define the velocity and chemical potential as
\begin{equation}\label{eq:mu_n}
v_n=\IM\left(\frac{\d_x\psi_n}{\psi_n}\right),\quad \mu_n=-\frac12\RE\left(\frac{\d_x^2\psi_n}{\psi_n}\right)+f'(|\psi_{n-1}|^2)+V_{n-1}.
\end{equation}
and we will see below this definition is consistent by the positivity of $|\psi_n|$. By \eqref{eq:QLS_n} the equation of $v_n$ is given by
\begin{equation}\label{eq:dtv_n}
\d_t v_n=-\d_x\left(\mu_n+\frac{1}{\tau}S_{n-1}\right).
\end{equation}
Thus following the idea of \eqref{eq:gradS}, we define $S_n$ to be the solution to the space-time gradient equation 
\begin{equation}\label{eq:gradS_n}
\begin{cases}
\d_xS_n =v_n,\quad\d_t S_n=-\mu_n-\frac{1}{\tau}S_{n-1} \\
S_n(0,0)=S_*\in [0,2\pi).
\end{cases}
\end{equation}
The "irrotationality" condition \eqref{eq:dtv_n} guarantees the solvability of \eqref{eq:gradS_n}, and following the idea of Definition \ref{def:phase}, the phase function $S_n$ can be explicitly written as
\begin{equation}\label{eq:defS_n}
\begin{aligned}
S_n(t,x)=&S_*+\int_\T\int_0^{x_*}v_0(y)dydx_*\\
&-\int_0^t \int_\T(\mu_n+\tau^{-1}S_{n-1})(s,x_*)dx_*ds+\int_\T\int_{x_*}^x v_n(t,y)dydx_*.
\end{aligned}
\end{equation}

We claim that the sequence $\{\psi_n\}$ satisfies the following properties on $[0,T_*)\times \T$:
\begin{itemize}
\item[(1)] boundedness: $\|\psi_n\|_{L^\infty_tH^2_x}\leq 2\|\psi_0\|_{H^2_x}$ and $\|S_{n}\|_{L^\infty_tH^2_x}\leq C$;
\item[(2)] positivity: $\inf_{t,x}|\psi_n|\geq \frac{\delta^\frac12}{2}$.
\end{itemize}
Here $0<C<\infty$ may depend on $\delta$ and $\|\psi_0\|_{H^2_x}$, but is independent of $n$. By our assumption, $\psi_0$ and $S_0$ obviously satisfy the properties (1) and (2). In this paper we will only focus on the estimates related to the phase function $S_n$ and the positivity property (2), and the remaining analysis following the standard argument for non-linear Schr\"odinger equation.

For the boundedness of $\|\psi_n\|_{L^\infty_tH^2_x}$, we use the Duhamel formula \eqref{eq:duhamel} and consider its $L^\infty_tH^2_x$ norm. The last integral in \eqref{eq:duhamel} can be estimated as
\begin{align*}
\|\int_0^te^{\frac{i}{2}(t-s)\d_x^2}S_{n-1}\psi_n(s)ds\|_{L^\infty_tH^2_x}\leq&  T_* \|S_{n-1}\psi_n \|_{L^\infty_tH^2_x}\\
\leq & T_* \|S_{n-1}\|_{L^\infty_tH^2_x}\|\psi_n \|_{L^\infty_tH^2_x}.
\end{align*}
Thus by the assumption of induction, we have
\[
\frac{1}{\tau}\|\int_0^te^{\frac{i}{2}(t-s)\d_x^2}S_{n-1}\psi_n(s)ds\|_{L^\infty_tH^2_x}\leq \frac{CT_*}{\tau} \|\psi_n \|_{L^\infty_tH^2_x},
\]
and we choose $T_*$ small such that $\frac{CT_*}{\tau}\leq \frac14$. Similarly, by Sobolev inequalities, the second integral in the right hand side of \eqref{eq:duhamel} can be bounded by
\begin{align*}
\|\int_0^te^{\frac{i}{2}(t-s)\d_x^2}[f'(|\psi_{n-1}|^2)&+V_{n-1}]\psi_n(s)ds\|_{L^\infty_tH^2_x}\\
\leq & T_*\|f'(|\psi_{n-1}|^2)+V_{n-1}\|_{L^\infty_{t,x}}\|\psi_n\|_{L_t^\infty H^2_x}\\
\leq & CT_*\|\psi_n\|_{L_t^\infty H^2_x}\leq \frac14\|\psi_n\|_{L_t^\infty H^2_x},
\end{align*}
which implies
\[
\|\psi_n\|_{L_t^\infty H^2_x}\leq \|\psi_0\|_{H^2_x}+\frac12\|\psi_n\|_{L_t^\infty H^2_x},
\]
namely the first part of property (1). 

Now we prove the positivity property (2). Let us define the density $\rho_n=|\psi_n|^2$, then from \eqref{eq:QLS_n} we obtain the equation of $\rho_n$ as
\[
\d_t\rho_n+\d_x (\rho_n v_n)=0.
\]

As a consequence, the total mass (and the average density since we assume $|\T|=1$) $M_n(t)=\int_\T\rho_n dx\equiv M_0$ is conserved. By Poincar\'e inequality,
\begin{equation}\label{eq:poincare}
\|\rho_n-M_0\|_{L^\infty_{t,x}}\leq \frac12\|\d_x\rho_n\|_{L^\infty_tL^1_x}\leq M_0^\frac12\|\d_x\psi_n\|_{L^\infty_tL^2_x},
\end{equation}
therefore 
\[
\inf_{t,x}|\psi_n|^2=\inf_{t,x}\rho_n\geq \frac{\delta}{4}
\]
if we have
\begin{equation}\label{eq:bd_pc}
\|\d_x\psi_n\|_{L^\infty_tL^2_x}\leq M_0^\frac12-\frac{\delta}{4}M_0^{-\frac12}.
\end{equation}
Then we consider the norm $\|\d_x\psi_n\|_{L^\infty_tL^2_x}$. Multiplying \eqref{eq:QLS_n} by $2\d_x^2\bar\psi_n$ and taking the imaginary part, it follows that
\begin{align*}
\frac{d}{dt}\int_\T|\d_x\psi_n|^2dx=&2\int_\T[f'(|\psi_{n-1}|^2)+\tau^{-1}S_{n-1}+V_{n-1}]\IM(\psi_n\d_x^2\bar\psi_n)dx\\
\leq & 2\|[f'(|\psi_{n-1}|^2)+\tau^{-1}S_{n-1}+V_{n-1}]\|_{L^\infty_{t,x}}\|\psi_n\|_{L^\infty_{t,x}}\|\d_x^2\psi_n\|_{L^\infty_tL^2_x}.
\end{align*}
By using the previous estimates of $f'(|\psi_{n-1}|^2)$, $S_{n-1}$, $V_{n-1}$ and property (1), it gives
\[
\frac{d}{dt}\int_\T|\d_x\psi_n|^2dx\leq \frac{C}{\tau}.
\]
Last, by using the initial condition \eqref{eq:ini_1} we obtain
\[
\int_\T|\d_x\psi_n|^2(t)dx\leq \int_\T|\d_x\psi_0|^2+\frac{CT_*}{\tau}\leq (M_0^\frac12-\delta M_0^{-\frac12})^2+\frac{CT_*}{\tau},
\]
which allows us to choose $T_*$ small such that \eqref{eq:bd_pc} holds. Thus we prove property (2). 

Last, we prove the remaining part of property (1), namely $\|S_n\|_{L^\infty_tH^2_x}\leq C$. By $\d_xS_n=v_n$, we only need to control $\|S_n\|_{L^\infty_tL^2_x}$ and $\|\d_x v_n\|_{L^\infty_tL^2_x}$.
From \eqref{eq:defS_n} we obtain
\[
|S_n(t,x)|\leq S_*+ \int_\T|v_0|dy+\int_0^t\int_T(|\mu_n|+\tau^{-1}|S_{n-1}|)dx_*ds+\int_\T|v_n|dy.
\]
Recalling that $v_n=\IM(\frac{\d_x\psi_n}{\psi_n})$, by the bound of $\psi_n$ and (2) it implies
\[
\int_\T|v_n|dy\leq 2\delta^{-\frac12}\int_\T |\d_x\psi_n|dy\leq \frac{C}{4}.
\]
Same bound holds for $\int_\T|v_0|dy$. By the assumption of induction we have $\|S_{n-1}\|_{L^\infty_tL^2_x}\leq C$, therefore
\[
\int_0^t\int_T\tau^{-1}|S_{n-1}|dx_*ds\leq \frac{C}{4}
\]
if we choose $T_*\leq \frac{\tau}{4}$. Last, by \eqref{eq:mu_n} we have
\[
\int_0^t\int_\T|\mu_n|dx_*ds\leq 2\delta^{-\frac12}\int_0^t\int_\T|\d_x^2\psi_n|dx_*ds+\int_0^t\int_\T|f'(|\psi_{n-1}|^2)|dx_*ds\leq \frac{C}{4},
\]
where in the last inequality we use $f\in C^2([0,\infty))$ and 
\[
|f'(|\psi_{n-1}|^2)|\leq C(\|f\|_{C^2},\|\psi_{n-1}\|_{L^\infty_{t,x}})\leq C(\|f\|_{C^2},\|\psi_{n-1}\|_{L^\infty_{t}H^2_x}).
\]
Thus we prove $\|S_n\|_{L^\infty_{t,x}}\leq C$. On the other hand, again by the definition of $v_n$, we have
\[
\|\d_xv_n\|_{L_t^\infty L^2_x}\leq 2\delta^{-\frac12}\|\d_x^2\psi_n\|_{L_t^\infty L^2_x}+2\delta^{-\frac12}\|(\d_x\psi_n)^2\|_{L_t^\infty L^2_x},
\]
where by embedding inequality,
\[
\|(\d_x\psi_n)^2\|_{L_t^\infty L^2_x}\leq \|\psi_n\|_{L_t^\infty H^2_x}^2\leq 4\|\psi_0\|_{H^2_x}^2\leq C.
\]
Thus we finish the proof of property (1).

Now we need to show $\{(\psi_n,S_n)\}$ is a contractive sequence, which is proved by a analogous argument as property (1) by considering the difference of \eqref{eq:duhamel} and its $L^\infty_tH^2_x$ norm. Since the estimates are standard and almost repetitive, we will not give the details.

\begin{lem}\label{lem:contr}
There exists $0<q<1$, such that the sequence $\{(\psi_n,S_n)\}$ satisfies
\[
\|\psi_n-\psi_{n-1}\|_{L^\infty_tH^2_x}+\|S_{n}-S_{n-1}\|_{L^\infty_tH^2_x}\leq Cq^{n-1}.
\]
\end{lem}

Lemma \ref{lem:contr} imply $\{(\psi_n,S_n)\}$ converges strongly to $(\psi,S)$ in $L^\infty([0,T_*),H^2(\T))$. Moreover by the convergence and the uniform lower bound of $|\psi_n|$, we have
\[
v_n=\IM\left(\frac{\d_x\psi_n}{\psi_n}\right)\to v=\IM\left(\frac{\d_x\psi}{\psi}\right)
\]
\[
-\d_x^2 V_n=|\psi_n|^2-\mathcal{C}(x)\to-\d_x^2 V=|\psi|^2-\mathcal{C}(x)
\]
and
\[
\mu_n=-\frac12\RE\left(\frac{\d_x^2\psi_n}{\psi_n}\right)+f'(|\psi_{n-1}|^2)+V_n\to -\frac12\RE\left(\frac{\d_x^2\psi}{\psi}\right)+f'(|\psi|^2)+V=\mu.
\] 
Also we can rewrite \eqref{eq:defS_n} as
\begin{align*}
S_n(t,x)=&e^{-\frac{t}{\tau}}\left(S_*+\int_\T\int_0^{x_*}v_0(y)dydx_*\right)\\
&-\int_0^t \int_\T e^{\frac{s-t}{\tau}}[\mu_n+\tau^{-1}(S_{n-1}-S_n)](s,x_*)dx_*ds+\int_\T\int_{x_*}^x v(t,y)dydx_*,
\end{align*}
then by passing to the limit we see $S$ satisfies \eqref{eq:defS}. Thus we prove $(\psi,S)$ is a solution to the Cauchy problem \eqref{eq:cauchy_NLS} on $[0,T_*)\times\T$, and the continuity in time follows directly from the Duhamel formula.

The proof of uniqueness follows a standard argument for Schr\"odinger equation by considering the difference equation, similar to Lemma \ref{lem:contr} of contraction.
\end{proof}

\end{document}
